\section{Background and General Goals}
\label{sec:background}

\subsection{Background}

An autonomous AI agent is a software system in which a large language model serves as the central reasoning engine, directing the selection and execution of actions within a real computing environment \cite{wang2024survey}. Unlike a conventional language model that generates a single textual response to a user query, an agent operates in a loop: it reasons over the current state of a task, selects an appropriate action---such as reading a file, calling an API, or executing a shell command---observes the result, and updates its reasoning accordingly \cite{yao2022react}. This capacity for iterative, environment-aware decision-making is what differentiates an agent from a static inference system.

The standard mechanism for equipping agents with capabilities is \textit{tool use}, in which the developer assembles a list of available functions and their input and output schemas, then includes this list in the agent's system prompt at the beginning of every task \cite{schick2023toolformer}. The underlying language model is then expected to identify the appropriate function from this flat list and invoke it with the correct parameters. Research has demonstrated that this approach is effective for small, well-defined capability sets; however, performance degrades measurably when the number of available tools exceeds approximately twenty, because irrelevant function definitions occupy context window space that would otherwise be available for task reasoning \cite{patil2023gorilla}. This degradation is particularly severe for locally deployed small language models, which operate within context windows substantially smaller than those of frontier cloud systems and are more sensitive to prompt length \cite{shi2024context, gemma4_2026}.

An alternative approach, implemented in the OpenClaw autonomous agent framework \cite{openclaw_2025} and theorised in research on embodied agent policies \cite{wang2024executable}, structures capabilities as self-contained knowledge artefacts called \textit{agentic skills}. Each skill is represented as a structured \texttt{SKILL.md} file stored on the local file system, containing three mandatory components: a set of applicability conditions defining when the skill is appropriate to use; a numbered execution policy specifying each procedural step the agent must follow; and a set of termination criteria indicating when the task has been completed successfully or should be aborted. These files are not loaded into the agent's context simultaneously. Instead, they are disclosed in three stages: the agent begins with only a compact metadata index of skill names and brief descriptions; it fetches the full instructional content of a specific skill only when the current task demands it; and it accesses the associated executable scripts exclusively at the moment of action. This hierarchical loading pattern, termed \textit{Progressive Disclosure}, is designed to preserve reasoning capacity, enforce behavioural constraints, and reduce hallucinated invocations.

This study tests whether progressive skill disclosure delivers measurable improvements over flat tool use when both paradigms are applied to the same tasks, evaluated using identical language models served through Ollama \cite{ollama_2024} for local inference and the Gemini 3.1 Pro \cite{google_gemini_2025} and GPT-4o \cite{openai_gpt4o_2024} APIs for cloud-based comparisons. All orchestration and measurement logic is implemented in Python, without reliance on additional agent frameworks, ensuring that the experimental results reflect the disclosure mechanism itself rather than framework-specific optimisations.

\subsection{Statement of Objectives}

This thesis pursues three specific objectives:

\begin{enumerate}
    \item To empirically compare the \textit{task completion accuracy} and \textit{trajectory quality} of an autonomous agent operating under flat tool use against the same agent operating under progressive skill disclosure, across a benchmark of 40 tasks spanning three levels of compositional complexity.
    \item To quantify the \textit{token efficiency} and \textit{operational cost} of both paradigms, measuring the reduction in prompt token consumption and total inference expenditure achieved by the three-tier disclosure model relative to full upfront tool exposure.
    \item To determine the conditions---defined by skill library size and task complexity---under which progressive disclosure provides the greatest improvement, thereby establishing practical guidelines for practitioners deploying autonomous agents with large capability sets.
\end{enumerate}

\subsection{Significance of the Project}

Production autonomous agent deployments in enterprise document processing, automated research generation, and IT operations management are increasingly constrained by the cost and reliability of large-scale tool orchestration \cite{significant_ai_2025}. The architectural choice between flat tool exposure and hierarchical skill disclosure has direct consequences for token expenditure, hallucination frequency, and compliance with operational policies. Despite this practical importance, developers currently make architectural decisions without access to empirical, independently validated performance data. This thesis provides that evidence through a reproducible experimental protocol, contributing both to the emerging body of knowledge on agentic evaluation methodology \cite{liu2023agentbench} and to the practical guidance available to engineers designing autonomous systems with large or growing capability libraries.
